\begin{lemma}
\label{editorial-source-lemma-base-change-essentially-of-finite-type}
The class of ring maps which are essentially of finite type is
preserved by base change. Similarly for essentially of finite
presentation.
\end{lemma}

\begin{proof}
Retain the original map $R\to S$ and a map $\rho:R\to R'$.
Write $S=U^{-1}A$ through a specified $R$-algebra isomorphism,
where $A$ has the required finite type or finite presentation
property. Put $A'=A\otimes_R R'$ and let $U'$ be the image
of $U$ under $a\mapsto a\otimes1$.
There is an isomorphism of the actual base-changed rings
$$
\Phi:S\otimes_R R'\longrightarrow (U')^{-1}A',
\qquad
(a/u)\otimes r'\longmapsto
\frac{a\otimes r'}{u\otimes1}.
$$
To construct it, the ring map $A\to(U')^{-1}A'$ sends
$a$ to $(a\otimes1)/1$ and inverts $U$, so extends to $S$.
Together with the map $R'\to(U')^{-1}A'$,
$r'\mapsto(1\otimes r')/1$, it induces the displayed map;
the two maps agree on the original $R$ by tensor balancing.
Conversely $A'\to S\otimes_R R'$ sends
$a\otimes r'$ to $(a/1)\otimes r'$ and sends $u\otimes1$
to a unit with inverse $(1/u)\otimes1$.
It therefore induces a map $\Psi:(U')^{-1}A'\to S\otimes_R R'$.
Every numerator is a finite sum of pure tensors and every
denominator in $U'$ is the image of an original $u\in U$;
on those fractions its formula is
$$
\Psi\left(\frac{\sum_{i=1}^q a_i\otimes r_i'}{u\otimes1}\right)
=\sum_{i=1}^q(a_i/u)\otimes r_i'.
$$
The construction from the localization maps proves independence
of both tensor expressions and fraction representatives.
The composites fix each pure tensor $(a/u)\otimes r'$ and
each displayed fraction, respectively, so they are identities.

If $a_1,\ldots,a_t$ generate $A$ as an $R$-algebra, their
images $a_i\otimes1$ generate $A'$ as an $R'$-algebra:
apply the original polynomial expression for each $a$ in any
finite sum of pure tensors, retaining its coefficients through
$\rho$. Thus finite type is preserved.
For finite presentation, write
$A=R[X_1,\ldots,X_t]/(P_1,\ldots,P_m)$.
The coefficient map gives
$$
A'\cong R'[X_1,\ldots,X_t]/(\rho(P_1),\ldots,\rho(P_m)),
$$
where for the full original polynomial
$P_j=\sum_\nu r_{j\nu}X^\nu$ one has
$\rho(P_j)=\sum_\nu\rho(r_{j\nu})X^\nu$.
Indeed $R[X_1,\ldots,X_t]\otimes_R R'$ maps isomorphically
to $R'[X_1,\ldots,X_t]$ by
$rX^\nu\otimes r'\mapsto\rho(r)r'X^\nu$, with inverse
$r'X^\nu\mapsto X^\nu\otimes r'$.
Tensoring the quotient presentation is right exact, so the
kernel of the resulting surjection is generated by the images
of the original ideal $(P_1,\ldots,P_m)$.
A finite sum of its polynomial multiples maps to the same
finite sum of multiples of $\rho(P_j)$, and each listed
$\rho(P_j)$ occurs, proving the asserted ideal equality.
This is a finite presentation of $A'$.
In either case the proved isomorphism $S\otimes_R R'\cong
(U')^{-1}A'$ establishes the required essential property.
It remains valid when a denominator maps to zero, since both
localization constructions and their inverse maps then yield
the zero ring.
\end{proof}